\section{Introduction}

The gap between ``I read the transformer paper'' and ``I understand how a
modern LLM is actually built, trained, aligned, and served'' has widened
dramatically. A 2025-era system is not a transformer: it is latent-compressed
attention over a paged KV cache, sparse experts balanced without auxiliary
losses, extra prediction heads that double as speculative drafts, a
preference-optimization and verifiable-reward RL pipeline, and a continuous-%
batching server---each the subject of its own paper, none of them teachable in
isolation.

Existing educational resources each cover one leg
(\S\ref{sec:related}): pre-training \citep{karpathy2022nanogpt,
raschka2024build}, single-paper architecture \citep{rush2018annotated,
jayasiri2020labml}, deep RL \citep{achiam2018spinningup}, or ML-systems
foundations \citep{tinytorch2026}. A learner who wants the \emph{connected}
story---why mid-training exists, how MTP heads become a serving speedup, why
the KV cache dictates attention design---must stitch it together from a dozen
sources with incompatible codebases.

BabyWhale closes this gap with one design decision: \textbf{a single tiny
model's biography is the curriculum} (Figure~\ref{fig:lifecycle}). The same
checkpoint is born (architecture), learns to read (pre-training), specializes
(mid-training), learns to behave (SFT, DPO), learns to reason (GRPO with a
code sandbox), is compressed (FP4), goes to work (paged KV, speculative
decoding, continuous batching), and is judged (bits-per-byte, pass@1,
long-context retrieval)---all on a laptop, in minutes per stage.

\paragraph{Target audience and impact.} BabyWhale serves three groups:
(i)~\emph{self-learners and graduate students} who know basic deep learning
and want the modern LLM stack end-to-end rather than in slices;
(ii)~\emph{instructors}, since each module is self-contained and its labs
autograde with a plain \texttt{python} invocation---no grading
infrastructure; and (iii)~\emph{research engineers onboarding onto LLM
systems}, for whom the measured ablations double as a systems-intuition
reference. Because everything runs in classroom time on a laptop, the
barrier to adoption is one machine, not a cluster.

Our second contribution addresses a failure mode we measured in ourselves:
\textbf{educational prose drifts from code}. While auditing our own drafts
against the implementation we found and fixed eight confidently wrong
explanations (a mis-stated gating function, a wrong loss attribution, an
inverted masking claim). BabyWhale therefore treats pedagogy as a tested
artifact: labs are graded against the real modules, pasted snippets are
CI-checked against the source, and every quantitative claim is a command the
reader can run (\S\ref{sec:course}, \S\ref{sec:measured}).

\begin{figure*}[t]
  \centering
  % Figure 1: lifecycle spine + module anatomy (TikZ)
\begin{tikzpicture}[
  font=\small,
  stage/.style={draw, rounded corners=2pt, align=center, minimum height=9mm,
                minimum width=27mm, inner sep=2pt, fill=blue!6},
  anat/.style={draw, rounded corners=2pt, align=left, inner sep=4pt, fill=gray!8},
  arr/.style={-{Stealth[length=2mm]}, thick}
]
% --- top row -----------------------------------------------------------
\node[stage] (born)  {\textbf{born}\\ \scriptsize arch.\ 01--07};
\node[stage, right=4mm of born]  (read)  {\textbf{learns to read}\\ \scriptsize pre-train 08--09};
\node[stage, right=4mm of read]  (spec)  {\textbf{specializes}\\ \scriptsize mid-train 10};
\node[stage, right=4mm of spec]  (beh)   {\textbf{learns to behave}\\ \scriptsize SFT/DPO 11--12};
% --- bottom row (right to left) ----------------------------------------
\node[stage, below=5mm of beh]   (reas)  {\textbf{learns to reason}\\ \scriptsize RL 13};
\node[stage, left=4mm of reas]   (comp)  {\textbf{gets compressed}\\ \scriptsize quant.\ 18};
\node[stage, left=4mm of comp]   (work)  {\textbf{goes to work}\\ \scriptsize serving 14--17};
\node[stage, left=4mm of work]   (judge) {\textbf{is judged}\\ \scriptsize eval 19};
% --- arrows ------------------------------------------------------------
\draw[arr] (born) -- (read);
\draw[arr] (read) -- (spec);
\draw[arr] (spec) -- (beh);
\draw[arr] (beh)  -- (reas);
\draw[arr] (reas) -- (comp);
\draw[arr] (comp) -- (work);
\draw[arr] (work) -- (judge);
\node[left=2.5mm of born, align=center, font=\footnotesize]
  (ckpt) {\emph{one}\\ \emph{checkpoint}};
\draw[arr, dashed] (ckpt) -- (born);
% --- anatomy strip ------------------------------------------------------
\node[anat, below=5mm of judge.south west, anchor=north west, text width=46mm]
  (beats) {\textbf{5 beats / module}\\[1pt] \scriptsize
  wall $\to$ idea $\to$ \emph{theory-to-code} table $\to$ measured payoff $\to$ break it};
\node[anat, right=4mm of beats.north east, anchor=north west, text width=44mm]
  (tracks) {\textbf{3 tracks}\\[1pt] \scriptsize
  Read (real code shown) \\ Build (23 labs, graded vs.\ real modules) \\ Extend (open PRs)};
\node[anat, right=4mm of tracks.north east, anchor=north west, text width=46mm]
  (guards) {\textbf{CI integrity}\\[1pt] \scriptsize
  drift guard: shown code $\equiv$ source \\ milestones: real-task gates \\ 334 tests, Apple-Silicon CI};
\end{tikzpicture}

  \caption{\textbf{BabyWhale in one view.} \emph{Top:} the course is one
  model's lifecycle; each stage names its course modules. \emph{Bottom:} every
  module has the same anatomy---five narrative beats read through three
  lenses, three participation tracks, and two CI-enforced integrity
  mechanisms that keep the teaching aligned with the code.}
  \label{fig:lifecycle}
\end{figure*}
